% Generated from locale/tr/content/second-order-logic/syntax-and-semantics/terms-formulas.tex; do not hand-edit.


\olfileid[tr]{sol}{syn}{frm}

\olsection{Terimler ve \printtoken{P}{formula}}

Birinci dereceden mantıkta olduğu gibi, ikinci dereceden mantığın ifadeleri
de \emph{!!{variable}s}, \emph{!!{constant}s}, \emph{!!{predicate}s} ve
bazen \emph{!!{function}s} içeren temel bir söz dağarcığından kurulur.
Bunlardan, mantıksal bağlaçlar, niceleyiciler ve parantezlerle virgüller gibi
noktalama simgeleriyle birlikte \emph{terimler} ve \emph{!!{formula}s}
oluşturulur. Aradaki fark, ikinci dereceden mantığın nesne değişkenlerine ek
olarak bağıntı ve fonksiyon değişkenleri de içermesi ve bunlar üzerinde
nicelemeye izin vermesidir. Dolayısıyla ikinci dereceden mantığın mantıksal
simgeleri, birinci dereceden mantığın simgelerine ek olarak şunlardır:

\begin{enumerate}
\item Her arite~$n$ için !!{denumerable}s bir ikinci dereceden bağıntı
  !!{variable}s kümesi: $\Obj V_0^n$, $\Obj V_1^n$, $\Obj V_2^n$, \dots
\item !!{denumerable}s bir ikinci dereceden fonksiyon !!{variable}s kümesi:
  $\Obj u_0^n$, $\Obj u_1^n$, $\Obj u_2^n$, \dots
\end{enumerate}

Birinci dereceden $\Obj v_i$ değişkenleri için üst dil değişkenleri olarak
$x$, $y$, $z$ kullandığımız gibi, $\Obj V_i^n$ için $X$, $Y$, $Z$, vb.;
~$\Obj u_i^n$ için de $u$, $v$, vb. kullanacağız.

\begin{explain}
İkinci dereceden bir dilin mantıksal olmayan simgeleri, birinci dereceden
bir dilinkilerle aynı biçimde, yani dilin !!{constant}s, !!{function}s ve
!!{predicate}s öğeleri listelenerek belirlenir.

Birinci dereceden mantıkta !!{identity}~$\eq$ genellikle dile katılır.
Birinci dereceden mantıkta bir~$\Lang{L}$ dilinin mantıksal olmayan
simgeleri ilginç herhangi bir şeyi ifade edebilmemiz için çok önemlidir.
Elbette mantıksal olmayan hiçbir simge kullanmayan !!{sentence}s vardır;
fakat yalnızca~$\eq$ ile ilginç bir şey söylemek güçtür. İkinci dereceden
mantıkta sınırsız sayıda bağıntı ve fonksiyon değişkenimiz olduğundan,
mantıksal olmayan simgelerden oluşan özel bir dağarcık olmadan bile birinci
dereceden bir dilde söyleyebileceğimiz her şeyi söyleyebiliriz.
\end{explain}

\begin{defn}[İkinci dereceden terimler]
$\Lang L$ dilinin \emph{ikinci dereceden terimler} kümesi $\TrmSOL[L]$,
\olref[fol][syn][frm]{defn:terms} tanımına şu madde eklenerek tanımlanır:
\begin{enumerate}
\item $u$, $n$-li bir fonksiyon değişkeni ve $t_1$, \dots, $t_n$ terimler
  ise $\Atom{u}{t_1, \ldots, t_n}$ bir terimdir.
\end{enumerate}
\end{defn}

\begin{explain}
Dolayısıyla ikinci dereceden bir terim, birinci dereceden bir terime tıpatıp
benzer; tek fark, birinci dereceden terimde !!a{function}~$\Obj{f^n_i}$
bulunan yerde ikinci dereceden terimde bir fonksiyon değişkeni~$\Obj{u^n_i}$
bulunabilmesidir.
\end{explain}

\begin{defn}[İkinci dereceden \usetoken{s}{formula}]
$\Lang L$ dilinin \emph{ikinci dereceden !!{formula}s} kümesi
$\FrmSOL[L]$, \olref[fol][syn][frm]{defn:formulas} tanımına şu maddeler
eklenerek tanımlanır:
\begin{enumerate}
\item $X$, $n$-li bir yüklem değişkeni ve $t_1$, \dots, $t_n$,
  $\Lang L$ dilinin ikinci dereceden terimleri ise
  $\Atom{X}{t_1,\ldots, t_n}$ atomik bir !!{formula}{}dür.

\tagitem{prvAll}{$!A$ !!a{formula} ve $u$ bir fonksiyon değişkeni ise
  $\lforall[u][!A]$ !!a{formula}{}dür.}{}

\tagitem{prvAll}{$!A$ !!a{formula} ve $X$ bir yüklem değişkeni ise
  $\lforall[X][!A]$ !!a{formula}{}dür.}{}

\tagitem{prvEx}{$!A$ !!a{formula} ve $u$ bir fonksiyon değişkeni ise
  $\lexists[u][!A]$ !!a{formula}{}dür.}{}

\tagitem{prvEx}{$!A$ !!a{formula} ve $X$ bir yüklem değişkeni ise
  $\lexists[X][!A]$ !!a{formula}{}dür.}{}
\end{enumerate}
\end{defn}

